\documentclass[10pt,a4paper]{amsart}
\usepackage[margin=19mm]{geometry}
\usepackage{amsmath,amssymb,mathtools}
\usepackage[scr=rsfs,cal=euler]{mathalfa}
\usepackage{xfrac}
\usepackage[unicode,hidelinks,bookmarks=false]{hyperref}
\newcommand{\ie}{i.e.\ }
\newcommand{\cf}{cf.\ }
\newcommand{\resp}{respectively\ }
\newcommand{\Subsection}{\subsection}
\providecommand{\nobreakdash}{\nobreak\hbox{-}}
\setlength{\emergencystretch}{2em}
\multlinegap0pt
\usepackage{bm}
\usepackage{bbm}
\usepackage{dsfont}
\usepackage{xspace}
\usepackage{cancel}
\usepackage{mathtools}
\usepackage{amsthm}
\makeatletter
\@ifundefined{theorem}{\newtheorem{theorem}{Theorem}}{}
\makeatother
\makeatletter
\DeclareMathOperator{\As}{As}
\newcommand{\CC}{\mathbf{C}}
\newcommand{\QQ}{\mathbf{Q}}
\newcommand{\Qp}{\QQ_p}
\newcommand{\ah}{\mathrm{ah}}
\newcommand{\cE}{\mathcal{E}}
\newcommand{\cF}{\mathcal{F}}
\newcommand{\cL}{\mathcal{L}}
\expandafter\ifx\csname smatrix\endcsname\relax
\newenvironment{smatrix}{\left( \begin{smallmatrix} } {\end{smallmatrix} \right) }
\fi
\newcommand{\uPi}{\underline{\Pi}}
\makeatother
\title[Poles of p-adic Asai L-functions and distinguished represent]{Poles of p-adic Asai L-functions and distinguished representations}
\author{David Loeffler, Sarah Livia Zerbes}
\date{Source: arXiv:2501.11472v2 (CC BY 4.0)}
\begin{document}
\maketitle

\noindent\emph{Source and licence.} Poles of p-adic Asai L-functions and distinguished representations. Version arXiv:2501.11472v2 (CC BY 4.0), distributed under the Creative Commons Attribution 4.0 International licence, as recorded in the arXiv licence metadata for this exact version and shown on the abstract page https://arxiv.org/abs/2501.11472v2. Full rights notice: licenses/arxiv-2501.11472-CC-BY-4.0.txt.\par\smallskip

\noindent\textit{Editorial scope.} Counted unit: the Theorem whose source label is \texttt{thm:nopole} in the article (source0.tex lines 498--505, source label \texttt{thm:nopole}), delivered together with the article's own complete original proof of it (source0.tex lines 507--525, one \texttt{proof} environment of 19 lines). No mathematics is written, repaired or re-derived: statement and proof are the article's own text line for line. Required context retained uncounted: none. Every definition, display and label of those lines is retained unchanged. Omitted from this package and not redistributed: 1--497, 506--506, 526--847. Nothing inside a retained excerpt is omitted or altered: every retained body is delivered exactly as the article's lines are written, so no omission receipt of any kind accompanies this package. Citations: the retained excerpts carry no citation command at all, so no bibliography entry of the article is transcribed and no reference is redistributed. Reference adaptation: none was needed and none was made; every reference inside the delivered unit resolves inside it, so no unresolved reference or citation is produced. Printed slips kept literal: no printed word, symbol, hypothesis or number of the counted statement or of its proof was altered. Layout and class: the article's own class is \texttt{amsart} with packages including amsmath, amscd, amssymb, mathrsfs, mathtools, fullpage, enumerate, thmtools, fancyhdr, stmaryrd, tikz-cd, etoolbox, chngcntr, yfonts; the wrapper is the lane's pinned \texttt{amsart} editorial wrapper at 10pt on A4, which restates the theorem-like environments the retained text needs and adds the article's own macro definitions that the retained text needs (\texttt{\textbackslash As}, \texttt{\textbackslash CC}, \texttt{\textbackslash QQ}, \texttt{\textbackslash Qp}, \texttt{\textbackslash ah}, \texttt{\textbackslash cE}, \texttt{\textbackslash cF}, \texttt{\textbackslash cL}, \texttt{\textbackslash uPi}), with extra wrapper packages (bm, bbm, dsfont, xspace, cancel, mathtools). The delivered numbering is the unit's own. Assets: no retained span contains a figure environment, a graphics-inclusion command, a PDF-page inclusion, a table or a TikZ picture; no image file is copied into this package, the package declares no asset dependency and no image file is redistributed with it. Licence: version arXiv:2501.11472v2 of this article is distributed under the Creative Commons Attribution 4.0 International licence, as recorded in the arXiv licence metadata for this exact version and shown on the abstract page https://arxiv.org/abs/2501.11472v2. A full rights notice is delivered at licenses/arxiv-2501.11472-CC-BY-4.0.txt. Characters: every retained excerpt is pure ASCII, so the delivered engine, XeLaTeX with its default Latin Modern text font, reports no missing character.
 \begin{theorem}
  \label{thm:nopole}
  The following are equivalent:
  \begin{itemize}
   \item $\cL^\flat_{p,\As}(\uPi)$ has a pole at $\kappa = 0$.
   \item $\Pi$ is distinguished.
  \end{itemize}
 \end{theorem}

 \begin{proof}
  Since  $\cL^\flat_{p,\As}(\kappa)$ can have at worst a simple pole at $\kappa=0$, the statement will follow if we can show that
  \[ \left(\ell(\kappa) \cdot \cL^\flat_{p,\As}(\kappa)\right)|_{\kappa=0}=0\]
  unless $\kappa$ is distinguished. However, this limit is by construction equal to
  \[
   \tfrac{p + 1}{p} \cdot (\sqrt{D})^{-k} \left\langle \iota^\star\left(\nu_{\underline{\Pi}}\middle|_{\kappa = 0}\right),\,
   \lim_{\kappa \to 0} \ell(\kappa) \cdot \cE^{\flat}_{1/N}(2\kappa) \right\rangle, \]
  where $\ell(-)$ is the ``logarithm'' function on weight space (which maps a character $\kappa$ to $\kappa'(1)$, so it takes the value $k$ at the character $x \mapsto x^k$; in particular it has a simple zero at the trivial character).
  
  
  we are abusing notation slightly by treating $\kappa$ as a 
  The vector $\nu_\Pi \coloneqq \nu_{\underline{\Pi}}|_{\kappa = 0}$ is a $\Qp$-basis of the 1-dimensional $\Pi$-eigenspace in the coherent $H^1$ of the Hilbert modular surface (with a suitable coefficient system depending on $k$). This is canonically the base extension to $\Qp$ of a one-dimensional $\QQ$-vector space; and the base-extension to $\CC$ of the same space is spanned by $\cF^{\ah, 1}$. Hence we can find scalars $\Omega_p \in \Qp^\times$ and $\Omega_\infty \in \CC^\times$ such that
  \[
   \frac{\nu_\Pi}{\Omega_p} = \frac{\cF^{\mathrm{ah}, 1}}{\Omega_\infty} \in H^1(X_1(\mathfrak{N}) / \QQ, \omega^{(2-k, k)})[\Pi].
  \]
  On the other hand, the limit on the right-hand side of the pairing is the non-zero constant $\tfrac{1}{2}\left(1 - \tfrac{1}{p}\right)$, considered as a weight 0 $p$-adic modular form. Thus we have shown
  \[ \frac{\left(\ell(\kappa) \cdot \cL^\flat_{p,\As}\right)|_{\kappa=0}}{\Omega_p(\Pi)} = \tfrac{1}{2} (\sqrt{D})^{-k} \left(1 - \tfrac{1}{p^2}\right)\cdot \frac{\mathcal{I}(\Pi)}{\Omega_\infty(\Pi)}. \]
  Since $\mathcal{I}(\Pi)$ is non-zero if and only if $\Pi$ is distinguished, the proof is complete.
 \end{proof}

\end{document}
